%----------------------------------------------------------------------------
\appendix
%----------------------------------------------------------------------------
\chapter*{\fuggelek}\addcontentsline{toc}{chapter}{\fuggelek}
\setcounter{chapter}{\appendixnumber}
%\setcounter{equation}{0} % a fofejezet-szamlalo az angol ABC 6. betuje (F) lesz
\numberwithin{equation}{section}
\numberwithin{figure}{section}
\numberwithin{lstlisting}{section}
%\numberwithin{tabular}{section}

\input{include/\AIdeclaration}

%----------------------------------------------------------------------------
\newpage
\section{A laborkörnyezet forrásfájljai}
%----------------------------------------------------------------------------

A melléklet a \texttt{pqc-lab-forras.zip} tömörített archívumban érhető el, amely a \texttt{pqc-lab/} Docker-alapú laborkörnyezet teljes forráskódját tartalmazza -- a tanúsítványokat (\texttt{certs/}) és a mérési eredményeket (\texttt{experiment-results/}) kivéve. Az archívum könyvtárszerkezete a következő:

\begin{lstlisting}[
    float=false,
    basicstyle=\small\ttfamily,
    language={}
]
pqc-lab/
|-- .env
|-- docker-compose.yml
|-- README.md
|-- base-image/
|   `-- Dockerfile
|-- pqc-proxy/
|   |-- Dockerfile
|   |-- entrypoint.sh
|   `-- stunnel.conf.template
|-- mosquitto-oqs/
|   |-- Dockerfile
|   |-- entrypoint.sh
|   |-- mosquitto.conf
|   `-- stunnel-server.conf.template
|-- iot-node/
|   |-- Dockerfile
|   |-- node.py
|   `-- requirements.txt
|-- prometheus/
|   `-- prometheus.yml
|-- grafana/
|   `-- provisioning/
|       |-- dashboards/
|       |   |-- dashboard.yml
|       |   `-- pqc-overview.json
|       `-- datasources/
|           `-- prometheus.yml
`-- scripts/
    |-- analyze_results.py
    |-- export_metrics.py
    |-- generate_certs.sh
    `-- run_experiment.sh
\end{lstlisting}

Az alábbiakban röviden ismertetem az egyes könyvtárak tartalmát.

%----------------------------------------------------------------------------
\subsection*{\texttt{.env} és \texttt{docker-compose.yml}}
%----------------------------------------------------------------------------

A \texttt{.env} fájl a kísérlet futtatható paramétereit tartalmazza alapértelmezett értékekkel: a vizsgált PQC-algoritmust (\texttt{PQC\_ALG}), a CPU-korlátot (\texttt{CPU\_LIMIT}), az MQTT publish-intervallumot és az üzenetméretet. A \texttt{docker-compose.yml} írja le a teljes laborkörnyezet service-eit: a PQC alap image-et, a három IoT-csomópontot, a PQC proxyt, az MQTT brókert, valamint a Prometheus és Grafana megfigyelhetőségi stacket, beleértve a cAdvisor és node-exporter komponenseket.

%----------------------------------------------------------------------------
\subsection*{\texttt{base-image/}}
%----------------------------------------------------------------------------

Az egyetlen \texttt{Dockerfile} Ubuntu~24.04 alapon fordítja le és telepíti a \texttt{liboqs} kvantumbiztos kriptográfiai könyvtárat, valamint az \texttt{oqs-provider} OpenSSL~3 providert. Ez az image szolgál minden PQC-képességet igénylő konténer (proxy, bróker) kiindulópontjaként.

%----------------------------------------------------------------------------
\subsection*{\texttt{pqc-proxy/}}
%----------------------------------------------------------------------------

A PQC proxy konténer a stunnel TLS-alagút kliens oldalát valósítja meg. A \texttt{stunnel.conf.template} sablonból az \texttt{entrypoint.sh} indítószkript állítja elő a futásidejű konfigurációt: behelyettesíti a \texttt{PQC\_ALG} környezeti változóból az aktuális algoritmus nevét a \texttt{curves} direktívába, elvégzi az OpenSSL TLS kézfogás-késleltetés-mérést, majd elindítja a stunnelt.

%----------------------------------------------------------------------------
\subsection*{\texttt{mosquitto-oqs/}}
%----------------------------------------------------------------------------

Az Eclipse Mosquitto MQTT bróker konténere, amelybe a stunnel TLS-alagút szerver oldala is integrálva van. A \texttt{mosquitto.conf} a bróker hálózati és hitelesítési paramétereit adja meg, a \texttt{stunnel-server.conf.template} pedig a TLS-terminációhoz szükséges sablont tartalmazza. Az \texttt{entrypoint.sh} generálja a végleges stunnel konfigurációt, majd párhuzamosan indítja el a stunnelt és a Mosquitto brókert.

%----------------------------------------------------------------------------
\subsection*{\texttt{iot-node/}}
%----------------------------------------------------------------------------

A szimulált IoT szenzorcsomópont Python-alapú implementációja. A \texttt{node.py} MQTT~QoS~1 üzeneteket publikál a proxyn keresztül, és Prometheus-kompatibilis metrikákat exportál: publish-ciklus késleltetési hisztogramot, PUBACK visszaigazolási arányt, újracsatlakozásszámot és RSS-memóriahasználatot. A \texttt{requirements.txt} felsorolja a szükséges Python-csomagokat (\texttt{paho-mqtt}, \texttt{prometheus-client}, \texttt{psutil}).

%----------------------------------------------------------------------------
\subsection*{\texttt{prometheus/} és \texttt{grafana/}}
%----------------------------------------------------------------------------

A \texttt{prometheus/prometheus.yml} konfigurálja a Prometheus gyűjtési célpontjait (IoT-csomópontok, cAdvisor, node-exporter). A \texttt{grafana/provisioning/} könyvtár tartalmazza a Grafana automatikus adatforrás-konfigurációját (\texttt{datasources/prometheus.yml}) és az előre elkészített PQC-áttekintő dashboardot (\texttt{dashboards/pqc-overview.json}).

%----------------------------------------------------------------------------
\subsection*{\texttt{scripts/}}
%----------------------------------------------------------------------------

\begin{itemize}
    \item \texttt{generate\_certs.sh} -- ECDSA P-256 tanúsítványlánc generálása (CA + szerver) az stunnel TLS-kapcsolathoz.
    \item \texttt{run\_experiment.sh} -- az összes algoritmus $\times$ CPU-szint kombináció automatikus, szekvenciális futtatása; minden kísérlet előtt újraindul a labor.
    \item \texttt{export\_metrics.py} -- a Prometheus HTTP API-n keresztül lekérdezi és fájlba menti az összes releváns metrikát a kísérlet végeztével.
    \item \texttt{analyze\_results.py} -- offline eredményelemzés és diagramgenerálás a mentett metrikaexportból.
\end{itemize}
